The GPEG range assembly instruction is a complicated beast.
It has five parameters in bytecode, and a full version of the
instruction exists in assembly, supporting each of these.
However, for readability, there also exist shorter versions of the instruction.
Therefore we say that the range assembly instruction is polymorph.
It has three possible forms:

Single char-like matching. This is the result of case-sensitive
(or: non-case-insensitive) matching individual bytes in strings or
hex literals. In Lua PEG, this is equivalent to the 'char' instruction.
The assembly below strictly matches the ASCII letter 'a'.
The first form is the short form. The second form is the full form.

\begin{myquote}
\begin{verbatim}
range 61
range false 8 61 61 ff
\end{verbatim}
\end{myquote}

The second form of the range instruction assembly, is a proper range.
These are the result of sets in grammar, or macroes.
The assembly below strictly matches the range from ASCII 'a' through 'z',
inclusive. Again, the first form is the short form, while the second
form is the full form.

\begin{myquote}
\begin{verbatim}
range 61 7a
range true 8 61 7a ff
\end{verbatim}
\end{myquote}

The third form of the range instruction assembly is the form to which
it is compiled when encountering anything else, such as:

\begin{itemize}
\item Alternatives.
\item Bitmasks, even for matching less than a full byte.
\end{itemize}

In this case, the range instruction will take the following parameters,
in order: truerange, nbits, from, until, mask. In this example, both
instructions are full form, where the first example represents an
alternative (lower and upper case 'a', for example the result of a
case insensitive string match). The second example represents a bitmask
of six bits.

\begin{myquote}
\begin{verbatim}
range false 8 41 61 ff
range false 6 38 10 10
\end{verbatim}
\end{myquote}
